\chapter{Conclusion and Future Work}
\label{ch:conclusion}

\section{Summary}

This thesis proposed a communication-efficient Hierarchical Federated Learning framework for traffic flow prediction in Intelligent Transportation Systems. The framework addresses two central challenges in federated traffic forecasting: the statistical heterogeneity caused by non-IID traffic sensor distributions, and the communication overhead associated with repeated model exchange in bandwidth-constrained edge networks.

The proposed framework combines DTW-based sensor clustering, quality-weighted hierarchical aggregation, and adaptive reliability-based client selection. Together, these components aim to reduce intra-cluster heterogeneity, enable controlled inter-cluster knowledge sharing, and lower communication overhead under dynamic client availability.

The framework was evaluated on the METR-LA dataset (207 sensors) through a seven-configuration cumulative ablation study using a lightweight GRU Sequence-to-Sequence model. Multi-seed evaluation (three independent seeds) of the baseline and final configurations provided strong evidence for the robustness of the results: the final configuration achieved a mean MAE of $3.620 \pm 0.049$~mph across seeds, compared with $3.837 \pm 0.015$ for Global FedAvg, a 5.7\% improvement with non-overlapping observed seed ranges. The total communication cost was 25.8~GB over 60 communication rounds. Cross-dataset validation on PEMS-BAY (325 sensors) achieved an MAE of 2.052~mph without dataset-specific architectural tuning.

\section{Answering the Research Questions}

The experimental results presented in Chapter~\ref{ch:results} provide answers to the research questions defined in Chapter~\ref{chap:intro}.

\begin{enumerate}
    \item \textbf{Does DTW-based clustering improve prediction accuracy over global FedAvg under non-IID traffic sensor distributions?}

    Yes. The transition from global FedAvg (Run~A, MAE = 3.832) to DTW-clustered models (Run~B, MAE = 3.753) yields a 2.1\% MAE improvement, the largest individual gain in the ablation study. Multi-seed evaluation further confirms this finding: across three seeds, the proposed framework achieves a mean MAE of $3.620 \pm 0.049$ compared with $3.837 \pm 0.015$ for Global FedAvg, with non-overlapping observed seed ranges. This supports the conclusion that grouping sensors by temporal pattern similarity reduces intra-cluster non-IID divergence and allows cluster models to specialize in similar traffic behaviors.

    \item \textbf{Does hierarchical knowledge sharing between clusters improve accuracy over isolated cluster training?}

    Yes, although the measured improvement reflects the combined effect of quality-weighted intra-cluster aggregation and hierarchical inter-cluster synchronization. Moving from isolated DTW-clustered training (Run~B, MAE = 3.753) to the quality-weighted hierarchical configuration (Run~C, MAE = 3.722) yields a 0.8\% MAE improvement. This suggests that combining quality-aware intra-cluster updates with controlled inter-cluster knowledge sharing improves performance while preserving cluster specialization.

    \item \textbf{Can adaptive client selection reduce communication cost without significantly degrading prediction performance?}

    Yes. Adaptive reliability-based client selection (Run~D) reduces total communication by 12.2\% (24.5 $\to$ 21.5~GB) at the cost of only a 0.4\% MAE increase (3.722 $\to$ 3.736). This trade-off is favorable in bandwidth-constrained edge deployments where communication efficiency is a primary objective.

   \item \textbf{Is traffic pattern similarity, measured using DTW, a more effective basis for inter-cluster aggregation than geographic proximity?}

    The results provide limited but positive evidence for this claim. Replacing geographic-distance-based inter-cluster weights with DTW-distance-based weights (Run~E $\to$ Run~F) yields a 0.1\% MAE improvement (3.700 $\to$ 3.696). While small in absolute terms and within single-seed random variation, this change is conceptually consistent with the DTW-based clustering strategy: clusters with similar temporal patterns benefit more from shared model knowledge, regardless of their physical distance.
\end{enumerate}

Taken together, the results provide an affirmative answer to the central research question: \textit{a temporal-pattern-aware hierarchical federated learning framework can improve traffic prediction accuracy while reducing communication cost under non-IID traffic sensor distributions}. Multi-seed evaluation demonstrates a robust 5.7\% mean MAE improvement over the Global FedAvg baseline, with non-overlapping observed seed ranges across three independent seeds. The framework also reduces communication by 29.1\% compared with CMBA-FL, the most directly comparable published federated baseline.

\section{Concluding Remarks}

Overall, this thesis provides evidence that temporal-pattern-aware clustering and lightweight hierarchical aggregation can improve federated traffic speed forecasting under non-IID sensor distributions. The proposed framework achieves robust improvement over Global FedAvg across three seeds and provides a favorable communication--accuracy trade-off compared with recurrent federated baselines. At the same time, the results show that graph-based centralized models still hold an accuracy advantage, motivating future work on privacy-preserving spatial modeling.

\section{Future Work}

Several directions for future research emerge from the limitations identified in Section~\ref{sec:discussion}.

\begin{enumerate}
    \item \textbf{Incorporating spatial information without centralized graph access.} The proposed framework deliberately avoids graph structure to avoid centralizing raw sensor data and to reduce dependence on global network metadata. However, the 14--20\% MAE gap between the proposed framework and centralized graph-based models such as DCRNN and Graph WaveNet suggests that spatial modeling could provide further improvements. Future work could explore privacy-preserving spatial mechanisms, such as federated graph construction from local adjacency information or proxy-node representations~\cite{zhou2026fedhint}, that recover spatial dependencies without requiring centralized access to the full graph adjacency matrix.

    \item \textbf{Dynamic cluster reassignment.} The current DTW-based clustering is performed once before training and remains fixed. In practice, traffic patterns evolve due to seasonal changes, road construction, and policy interventions. Future work could investigate periodic re-clustering during training, triggered by monitoring intra-cluster loss divergence. The key challenge is balancing the potential accuracy gain from adaptive clusters against the additional communication and computation cost of cluster reassignment.

    \item \textbf{Cluster-adaptive model architectures.} All experiments in this thesis use the same GRU Seq2Seq architecture across all clusters. However, the per-cluster analysis reveals substantial differences in predictive difficulty (MAE range from 2.000 to 5.177), suggesting that clusters with more complex traffic patterns could benefit from larger or more expressive models. Future work could explore cluster-specific model capacity, where smaller clusters with regular patterns use lightweight models and larger clusters with heterogeneous patterns use deeper architectures.

    \item \textbf{Real-world edge deployment and availability modeling.} Client availability in this thesis is simulated using predefined reliability profiles. Real ITS edge networks exhibit more complex behavior, including correlated outages during extreme weather, variable bandwidth under peak network load, and hardware-specific computation delays. Deploying the framework on a physical testbed or using realistic network emulation would provide stronger validation of the adaptive client selection mechanism.

    \item \textbf{Extended cross-dataset evaluation.} The PEMS-BAY experiments were limited to 50 communication rounds due to computational constraints. Future work should evaluate the framework on additional traffic datasets with fully optimized training budgets, including datasets with different sensor densities, road types, and geographic regions, to more rigorously establish generalizability.

    \item \textbf{Communication compression techniques.} The current framework transmits full model parameters at each communication round. Future work could revisit gradient compression techniques, such as Top-K sparsification~\cite{gao2024fedce} or quantization, in combination with warm-up training or adaptive compression schedules. This could further reduce upload communication while avoiding convergence degradation during early federated training.
\end{enumerate}
